\documentclass[11pt,a4paper]{article}
\usepackage[utf8]{inputenc}
\usepackage[T1]{fontenc}
\usepackage[spanish]{babel}
\usepackage[margin=2.5cm]{geometry}
\usepackage{graphicx}
\usepackage[hidelinks]{hyperref}
\usepackage{enumitem}
\usepackage{listings}
\usepackage{xcolor}

\definecolor{codigo}{RGB}{235,235,235}

\lstset{
  basicstyle=\ttfamily\small,
  backgroundcolor=\color{codigo},
  breaklines=true,
  frame=single
}

\title{Examen Unidad IV\\
  Aplicaciones Web}
\author{Tejada Bajaña Luis Alejandro}
\date{28 de Agosto de 2026}

\begin{document}

% ==================== CARATULA ====================
\begin{titlepage}
  \centering

  {\Large \textbf{UNIVERSIDAD TÉCNICA ESTATAL DE QUEVEDO}}\\[0.4cm]
  {\large Facultad de Ciencias de la Computación y Diseño Digital}\\[0.2cm]
  {\large Carrera de Ingeniería de Software (Rediseño)}\\[1.2cm]

  \rule{\textwidth}{0.5pt}\\[0.6cm]
  {\LARGE \textbf{Examen Teórico-Práctico --- Unidad IV}}\\[0.3cm]
  {\Large \textbf{Modelo Vista Controlador y Servicios Web:\newline
  API REST, Autenticación JWT y Consumo de APIs Externas}}\\[0.6cm]
  \rule{\textwidth}{0.5pt}\\[1.2cm]

  {\large \textbf{Asignatura:} Aplicaciones Web (Quinto nivel, Paralelo A)}\\[0.4cm]
  {\large \textbf{Docente:} Dr. Gleiston Cicerón Guerrero Ulloa, Ph.D.}\\[0.4cm]
  {\large \textbf{Periodo:} 2026--2027 PPA}\\[1.2cm]

  {\Large \textbf{Estudiante:} Tejada Bajaña Luis Alejandro}\\[0.4cm]
  {\large \textbf{Número de carnet:} 1207432939}\\[0.4cm]
  {\large \textbf{Correo institucional:} \texttt{ltejadab@uteq.edu.ec}}\\[0.4cm]
  {\large \textbf{Fecha:} 28 de Agosto de 2026}\\[1.2cm]

  \textbf{Repositorio (URL):}\\[0.2cm]
  {\large \texttt{\url{https://github.com/Alxjandr07/App-Web-segundo-corte-Tejada-Luis}}}\\[1cm]

  \vfill
\end{titlepage}

% ==================== SECCIONES ====================
\section{Introducción}
El presente informe documenta el desarrollo del examen de la Unidad IV de la
asignatura Aplicaciones Web. El objetivo de la unidad es exponer la funcionalidad
del catálogo bibliotecario como una \textbf{API REST} con autenticación por
\textbf{JWT} (sin estado), así como el consumo de una \textbf{API externa}
(Open Library) aplicando el patrón \emph{cache-aside} sobre Redis y el manejo
diferenciado de fallos.

Todo el código fuente, el cuestionario teórico y las pruebas de integración se
encuentran versionados en el repositorio indicado en la carátula.

\section{Objetivos}
\begin{itemize}
  \item Exponer el catálogo como API REST respetando las restricciones de
    Fielding y envolviendo las respuestas exitosas en \texttt{ApiResponse}.
  \item Proteger las escrituras con roles mediante JWT \emph{stateless}.
  \item Consumir la API externa de Open Library con caché y manejo de errores.
  \item Verificar la solución con pruebas de integración.
\end{itemize}

\section{Desarrollo de la Parte B}

\subsection{B1. Listado de libros con filtros y paginación}
En \texttt{LibroController} se implementó el endpoint
\texttt{GET /api/v1/libros}, que acepta los parámetros opcionales
\texttt{titulo}, \texttt{categoriaId} y \texttt{anioDesde}, los delega a
\texttt{LibroService.buscar(...)} y devuelve un \texttt{Page<Libro>} paginado.
La respuesta se envuelve en \texttt{ApiResponse\{success, data, message, errors,
meta\}}, donde \texttt{meta} proviene de \texttt{PageMeta}.

\subsection{B2. Alta protegida por rol}
El endpoint \texttt{POST /api/v1/libros} recibe un \texttt{LibroRequest}
validado con \texttt{@Valid}. Un cuerpo inválido produce \texttt{400} en
formato \emph{Problem Details} con el arreglo \texttt{errors} poblado. La
creación correcta devuelve \texttt{201 Created} con la cabecera
\texttt{Location}. La operación está protegida por \texttt{@PreAuthorize} para
que solo el rol \texttt{ADMIN} pueda invocarla; un usuario con otro rol recibe
\texttt{403} y una petición sin token recibe \texttt{401}.

\subsection{B3. Consumo de la API externa con caché}
\texttt{OpenLibraryClient.consultarPorIsbn(String isbn)} consume
\texttt{https://openlibrary.org/isbn/\{isbn\}.json} con el bean
\texttt{RestClient} ya configurado (con \emph{timeouts}). Aplica
\texttt{@Cacheable} sobre el espacio \texttt{openlibrary} (TTL de 24 horas).
Los fallos se tratan de forma diferenciada: un \texttt{404} del proveedor
devuelve campos externos nulos y el endpoint responde \texttt{200} con el libro
local; un \texttt{5xx}, un \texttt{4xx} distinto de \texttt{404} o un
\emph{timeout} lanzan \texttt{ServicioExternoException}, convertida por el
manejador global en \texttt{502}. Los fallos nunca se cachean.

\subsection{B4. Pruebas de integración}
Se escribieron pruebas de integración en \texttt{LibroControllerIT}, que
heredan de \texttt{BaseIntegracionTest} y levantan un PostgreSQL real en un
contenedor efímero con Testcontainers.

\section{Resultados}
\begin{itemize}
  \item \texttt{mvn -B clean package -DskipTests} termina en \texttt{BUILD SUCCESS}.
  \item \texttt{mvn -B test} deja las pruebas de integración en verde.
  \item La caché de la API externa se evidencia con
    \texttt{docker compose exec redis redis-cli KEYS "openlibrary*"}.
\end{itemize}

\section{Conclusiones}
Se expuso y protegió el catálogo bibliotecario como una API REST con
autenticación JWT sin estado, se consumió una API externa aplicando
\emph{cache-aside} con manejo diferenciado de fallos y se verificó el
cumplimiento de los objetivos con pruebas de integración. El código se
encuentra disponible en el repositorio de la carátula.

\begin{thebibliography}{9}
\bibitem{fielding} R. T. Fielding, \emph{Architectural Styles and the Design of
  Network-based Software Architectures}. Ph.D. dissertation, University of
  California, Irvine, 2000.
\bibitem{jwt} M. B. Jones, J. Bradley y S. Sakimura, \emph{JSON Web Token
  (JWT)}, RFC 7519, IETF, mayo 2015.
\bibitem{rfc9457} M. Nottingham, E. Wilde y S. Dalal, \emph{Problem Details for
  HTTP APIs}, RFC 9457, IETF, julio 2023 (obsoleta RFC 7807).
\bibitem{http} R. T. Fielding, M. Nottingham y J. Reschke, \emph{HTTP
  Semantics}, RFC 9110, IETF, junio 2022.
\end{thebibliography}

\end{document}
